% =============================================================================
% CONCLUSION
% =============================================================================

We evaluated interwhen-style procedural input verification on a
clinical tool-calling benchmark across nine conditions spanning three
groups: capability anchors (A, B, B$'$), a primary study of four
extractor pipelines under best-practice configuration (upfront,
reactive, reactive + citations, reactive + k-shot), and two
exploratory verifier-mechanism comparators (C, D). Six paired
contrasts were pre-registered at family $\alpha=0.05$.

Best-practice verifier-guided extraction did not move accuracy above
B$'$. All four extractor pipelines matched the forced-tool baseline
($71.6$--$72.5\%$) with no significant contrast, and no architectural
axis --- placement, citation grounding, or $k$-shot voting --- carried
an effect. System accuracy is set by forced tool use ($+21.3$ pp over
plain tool access); the input verifier, at best-practice configuration,
neither helps nor hurts. Among the pipelines, reactive extraction is
the deployment choice --- identical accuracy at $\sim$$40\%$ of
upfront's token cost. The exploratory alternatives are consistent:
prompting the model to self-check (C) is no better than B$'$, post-hoc
output review (D) is no better than B$'$, and the external verifier
beats prompt-only self-check head-to-head.

Two design choices in the primary study --- the schema-gated verifier
and the query-style intervention --- were motivated by pilot
diagnostics (\cref{sec:methods_pilots}) and address by construction two
failure modes those pilots exposed. That gap proves to be null:
best-practice hygiene removes the harm the un-gated pilot caused but
produces no benefit over the forced-tool baseline, locating the binding
constraint in the LLM-extracted reference rather than the verifier
mechanism itself.
